\begin{abstract}
Aerial scene classifiers have to tell apart land-use classes that look alike, and they are usually trained on collections where some classes are rare.
We revisit a standard comparison on the 15-class SkyView dataset: hand-crafted pipelines (local binary patterns, colour histograms and SIFT bag-of-words with $k$-nearest-neighbour, SVM and stacking classifiers) against fine-tuned ResNet-18 and EfficientNet-B3, plus voting ensembles of both.
Every model is tuned on the same validation set and scored on the same balanced test set, and we train on three sets that differ only in their class distribution: balanced, long-tailed (16:1) and rebalanced by offline augmentation.
With balanced training the CNNs reach 96.2\% and 96.5\% macro-F1, against 84.4\% for the best hand-crafted pipeline.
The long tail costs the CNNs 2.4--4.0 points and the hand-crafted pipelines 9--13; stacking recognises only 2.5\% of the rarest class.
Rebalancing by augmentation helps the hand-crafted pipelines but lowers both CNNs by 1.1--1.7 points, while weighting the loss by inverse class frequency raises them by 0.9--1.2.
Averaging the two CNNs beats the better network in every setting (98.0\% macro-F1 on balanced data, McNemar $p<10^{-4}$); adding hand-crafted members does not help.
Finally, we show that breaking voting ties by member order turns any two-member majority vote into a copy of its first member, which is why these ensembles previously appeared to collapse.
\end{abstract}

\begin{IEEEkeywords}
aerial scene classification, remote sensing, transfer learning, class imbalance, ensemble learning, Grad-CAM
\end{IEEEkeywords}
